% Generated draft: requires mathematical and source review.
\documentclass[11pt]{article}
\usepackage[margin=25mm]{geometry}
\usepackage{fontspec}
\setmainfont{Latin Modern Roman}
\usepackage{amsmath,amssymb,mathtools,unicode-math}
\setmathfont{Latin Modern Math}
\usepackage{enumitem,needspace,xurl,hyperref}
\setlength{\parindent}{0pt}
\setlength{\parskip}{5pt}
\setlength{\emergencystretch}{3em}
\newcommand{\N}{\mathbb{N}}
\newcommand{\Z}{\mathbb{Z}}
\newcommand{\Q}{\mathbb{Q}}
\newcommand{\R}{\mathbb{R}}
\newcommand{\C}{\mathbb{C}}
\newcommand{\F}{\mathbb{F}}
\newcommand{\A}{\mathbb{A}}
\newcommand{\PP}{\mathbb{P}}
\newcommand{\E}{\mathbb{E}}
\newcommand{\eps}{\varepsilon}
\newcommand{\dd}{\,\mathrm{d}}
\begin{document}
\section*{UnsolvedMath ID: 30002933}
\textbf{Source problem code:} catalogue entry
\subsection{Definitions and mathematical statement}
The supplied statement asks three isoperimetric questions for a space \(X\) and a prescribed volume \(v\), but does not specify the geometric structure of \(X\), the admissible regions, or the meanings of volume and boundary area. Consequently, the following is a conditional formulation of the questions, requiring those data to be supplied.

Let \(\mathcal A\subseteq\mathcal P(X)\) be a specified class of admissible regions, and let
\[
V:\mathcal A\longrightarrow[0,+\infty],\qquad
P:\mathcal A\longrightarrow[0,+\infty]
\]
be specified volume and boundary-area functionals. Here \(P\) may represent an appropriate perimeter, but no particular perimeter definition is asserted by the source. These functionals and their extended nonnegative ranges are formal placeholders; the intended domains and finiteness requirements remain unspecified.

For a prescribed value \(v\) in the intended volume range, define
\[
\mathcal A_v:=\{E\in\mathcal A:V(E)=v\},\qquad
I_X(v):=\inf_{E\in\mathcal A_v}P(E),
\]
using the explicit convention \(\inf\varnothing=+\infty\). Define the collection of minimizers by
\[
\mathcal M_X(v):=
\{E\in\mathcal A_v:P(E)=I_X(v)\}.
\]
The notation suppresses the dependence on \(\mathcal A,V,P\).

\textbf{Questions.} For the intended space and prescribed volume:
\begin{enumerate}
\item Determine the isoperimetric profile value \(I_X(v)\), the infimum of the boundary area needed to enclose volume \(v\).
\item Determine whether this infimum is attained, that is, whether \(\mathcal M_X(v)\ne\varnothing\).
\item When it is attained, describe or classify all regions in \(\mathcal M_X(v)\).
\end{enumerate}
The source does not specify which spaces or volumes are quantified over, whether only finite-volume and finite-perimeter regions are intended, whether boundary area is relative to an ambient domain, or which equivalence relation should be used in a classification. If \(I_X(v)=+\infty\), the displayed definition permits infinite-area attainment; whether this is intended to count as an optimal shape is unresolved. These questions therefore specify a general research framework, rather than a uniquely determined conjecture for arbitrary spaces.

The supplied dated assessment discusses important essentially non-branching metric-measure spaces with synthetic lower Ricci bounds and reports partial results there. It does not supply the precise hypotheses or conclusions of those results, and those restrictions are not imposed here as hypotheses of the original question.
\subsection{Short English statement}
For a specified space, notion of volume, and notion of boundary area, ask for the least boundary area enclosing a prescribed volume, whether a region attains that value, and how all such regions can be described. The source leaves the geometric setting and admissibility conditions unspecified, so this is a conditional statement requiring human mathematical review.
\subsection{Sources}
\begin{itemize}
\item Local registry: supplied UnsolvedMath source record, permanent ID \(30002933\), source problem code \texttt{catalogue entry}, title \emph{Isoperimetric Regions of Fixed Volume in General Spaces}, in the Oberwolfach Reports Open Problems collection. The project identifies \texttt{data/unsolvedmath/problems.json} and \texttt{catalog/problems.json} as local registries; the supplied record's exact location within them was not independently checked.
\item Supplied report citation: \emph{Partial Differential Equations (2016)}, Oberwolfach Reports, DOI \texttt{10.4171/owr/2015/36}, pp.\ 2065--2124. \url{https://doi.org/10.4171/owr/2015/36}. The record gives proposed year \(2015\) and names no individual proposer. The report-title year \(2016\) and the \(2015\) provenance fields are preserved without reconciliation.
\item Supplied literature reference: Cavalletti--Mondino, ``sharp and rigid isoperimetry.'' \url{https://arxiv.org/abs/1502.06465}.
\item Supplied literature reference: Cavalletti--Mondino, ``finite-perimeter isoperimetry.'' \url{https://arxiv.org/abs/1610.05044}.
\item Literature assessment dated \texttt{2026-08-21}: the supplied curation labels the problem \texttt{partially\_solved} and reports sharp isoperimetric inequalities, rigidity, and classification of optimizers for important essentially non-branching spaces with synthetic lower Ricci bounds, while stating that the three questions are not solved for an arbitrary metric-measure space \(X\). This is supplied dated curation derived from the OWR record, not an independently verified literature assessment; its status and cited claims require independent verification.
\item The registry labels the statement \texttt{exact} and says it was retained byte-for-byte because it is internally readable and complete. This is supplied registry verification metadata, not independent verification against the report.
\end{itemize}
\end{document}
